%=============================================================================
%  05_ingestion.tex  -  Chapter 4: Data Ingestion and Profiling
%=============================================================================
\chapter{Data Ingestion and Profiling}
\label{ch:ingestion}

The journey of every dataset begins here. Before the language model sees
anything, the backend must turn an arbitrary uploaded file into a clean
\code{pandas.DataFrame}, understand its structure, and describe it in a form that
is both machine-usable (for retrieval) and human-readable (for the profile
panel and the report).

\section{The ingestion engine}
The \code{IngestionEngine} accepts raw bytes plus a filename and dispatches on
the file extension. Five formats are supported, declared explicitly so that an
unsupported upload fails with a clear message rather than a stack trace.

\begin{table}[h]
\centering
\caption{Supported input formats and their parsing strategy.}
\begin{tabularx}{\linewidth}{@{}L{2.4cm} L{2.2cm} X@{}}
\toprule
\textbf{Format} & \textbf{Extension} & \textbf{Strategy} \\
\midrule
CSV & \code{.csv} & \code{pandas.read\_csv} over an in-memory buffer. \\
JSON & \code{.json} & Tolerant parse: arrays become rows, single objects are
normalized, NDJSON is a fallback. \\
Parquet & \code{.parquet} & \code{pandas.read\_parquet} (pyarrow). \\
PDF & \code{.pdf} & \code{pdfplumber} extracts tables page-by-page. \\
SQLite & \code{.sqlite}, \code{.db} & \code{sqlite3} reads all tables; the
largest is used. \\
\bottomrule
\end{tabularx}
\end{table}

\begin{lstlisting}[style=py,caption={Format dispatch in \code{IngestionEngine.parse\_file}.}]
SUPPORTED_EXTENSIONS = {"csv", "json", "parquet", "pdf", "sqlite", "db"}

class IngestionEngine:
    @staticmethod
    def parse_file(file_bytes: bytes, filename: str) -> pd.DataFrame:
        extension = filename.rsplit(".", 1)[-1].lower() if "." in filename else ""

        if extension == "csv":
            try:
                return pd.read_csv(io.BytesIO(file_bytes))
            except Exception as exc:
                raise ValueError(f"Could not read CSV file: {exc}") from exc
        elif extension == "json":
            return IngestionEngine._parse_json(file_bytes)
        elif extension == "parquet":
            return pd.read_parquet(io.BytesIO(file_bytes))
        elif extension == "pdf":
            return IngestionEngine._parse_pdf(file_bytes)
        elif extension in ("sqlite", "db"):
            return IngestionEngine._parse_sqlite(file_bytes)
        else:
            raise ValueError(
                f"Unsupported file format: '{extension}'. "
                f"Supported: {SUPPORTED_EXTENSIONS}")
\end{lstlisting}

\subsection{Tolerant JSON parsing}
JSON is the least predictable of the tabular formats. The engine accepts three
shapes: a top-level array of records (each object becomes a row via
\code{json\_normalize}), a single object (normalized to one row, or expanded when
its values are equal-length lists), and newline-delimited JSON (detected as a
fallback when a strict parse fails). A \code{utf-8-sig} decode transparently
strips a byte-order mark if present.

\subsection{PDF table extraction}
For PDF uploads the engine walks every page with \code{pdfplumber} and calls
\code{extract\_tables}. Header-only or empty tables are skipped; rows whose
column count does not match the header are filtered out. When multiple tables
share an identical column set they are concatenated; otherwise the largest table
is returned. This heuristic favours the single most substantial table in a
document rather than fragile multi-table stitching.

\subsection{SQLite handling}
Because \code{sqlite3} requires a file on disk, the engine writes the uploaded
bytes to a temporary file, enumerates user tables from \code{sqlite\_master},
reads each into a DataFrame, and returns the one with the most rows. The
temporary file is always removed in a \code{finally} block, and a corrupt or
non-database file yields a clear \code{ValueError} rather than a raw driver
error.

\begin{uawarn}[One table per upload]
Both the PDF and SQLite paths deliberately reduce a multi-table source to a
single DataFrame (the largest). This keeps the downstream profile, retrieval
index, and charts unambiguous. Analysing several tables from one file is noted
as future work.
\end{uawarn}

\section{The data profiler}
Once a DataFrame exists, the \code{DataProfiler} produces a
\code{DataProfile} --- a structured, JSON-safe description consumed by the
front-end overview panel, the RAG context, the chart planner, and the PDF
report. Profiling is entirely deterministic: no model is involved.

\subsection{Semantic typing}
Beyond the raw pandas dtype, the profiler assigns each column a \emph{semantic
type}, which is what actually drives chart selection and analysis. The
classification cascade is:

\begin{enumerate}
  \item \textbf{boolean} --- a true boolean dtype, or a two-value string column
        whose tokens are all boolean-like (\code{true/false}, \code{yes/no},
        \code{y/n}, \code{0/1}, \code{on/off}).
  \item \textbf{datetime} --- a native datetime dtype, or an object column that
        parses to a date for more than 80\% of a 200-row sample.
  \item \textbf{id} --- uniqueness above 0.9 combined with either an
        ``id''-bearing name or an integer-like dtype.
  \item \textbf{continuous} --- any other numeric column.
  \item \textbf{categorical} / \textbf{text} --- object columns split by
        cardinality against a \code{max(20, 5\% of rows)} threshold.
\end{enumerate}

\begin{lstlisting}[style=py,caption={Excerpt of the semantic-type cascade in \code{DataProfiler}.}]
if pd.api.types.is_bool_dtype(series):
    return "boolean"
if not pd.api.types.is_numeric_dtype(series) and DataProfiler._looks_boolean(series):
    return "boolean"
if pd.api.types.is_datetime64_any_dtype(series):
    return "datetime"
# ... object-that-parses -> datetime ...
uniqueness = (unique_count / row_count) if row_count else 0.0
name_has_id = "id" in str(name).lower()
integer_like = pd.api.types.is_integer_dtype(series)
if uniqueness > 0.9 and (name_has_id or integer_like):
    return "id"
if is_numeric:
    return "continuous"
\end{lstlisting}

\subsection{Per-column statistics}
For continuous and id columns the profiler records min, max, mean, median,
standard deviation, \textbf{skewness}, and an \textbf{IQR-based outlier count}
(values beyond $Q_1 - 1.5\,\mathrm{IQR}$ or $Q_3 + 1.5\,\mathrm{IQR}$). Datetime
columns contribute ISO-formatted min/max. Categorical, id, and boolean columns
record their top five values with counts. Every value passes through a
\code{\_json\_safe} helper that converts numpy and pandas scalars, timestamps,
and NaNs into plain JSON primitives --- essential because this profile is
serialized to the front end and embedded in prompts.

\begin{table}[h]
\centering
\caption{Fields captured per column in \code{ColumnProfile}.}
\small
\begin{tabularx}{\linewidth}{@{}L{4.2cm} X@{}}
\toprule
\textbf{Field} & \textbf{Meaning} \\
\midrule
\code{dtype}, \code{semantic\_type} & Raw pandas dtype and inferred role. \\
\code{null\_count}, \code{missing\_pct} & Absolute and percentage missingness. \\
\code{unique\_count} & Distinct non-null values. \\
\code{min/max/mean/median/std} & Summary statistics for numeric columns. \\
\code{skewness}, \code{outlier\_count} & Distribution shape and IQR outliers. \\
\code{sample\_values}, \code{top\_values} & Examples and most-frequent values. \\
\bottomrule
\end{tabularx}
\end{table}

\subsection{Dataset-level signals}
At the whole-table level the profiler computes overall missingness, an estimated
in-memory footprint (\code{memory\_usage(deep=True)}), and --- when at least two
continuous columns exist --- a rounded \textbf{correlation matrix}. It also
publishes convenience lists of the datetime, numeric, and categorical column
names, which the chart planner uses to decide which visualizations are even
possible.

\section{Automated EDA}
Layered on top of the profile, \code{generate\_auto\_eda} assembles a compact
exploratory summary --- human-readable size figures, de-duplicated column
groupings, and headline observations --- that seeds both the overview panel and
the opening context the model receives. Because it is derived from the
deterministic profile, the EDA carries the same trust guarantee: it reflects the
data, not the model's guesswork.

\begin{uanote}[Why deterministic profiling matters]
Everything the model later says about the data is anchored to this profile. By
computing types, statistics, and correlations \emph{without} the LLM, the system
guarantees a stable, reproducible factual base --- the model interprets, it does
not measure.
\end{uanote}
